% Integrated FF.42; review copy only, do not input separately.
\begin{claim}{FF.42}{Proved}\label{clm:FFmetabeliankernel}
Put \(F_p=K_p^{\rm ab}\), and define
\[
 W_p=\{a\in B_p^\times:
                  a^m\in F_p^\times
                  \text{ for some integer }m\ge1\}.
                                                        \tag{FF.42.1}
\]
For every finite extension \(E/B_p\),
let \(\kappa_E\) denote the rational
\(p\)-primary Kummer map. Then
\[
 \ker\bigl(\kappa_E:E^\times
             \longrightarrow H^1(G_E,\mathbb Q_p(1))\bigr)
       =\bigcap_{n\ge1}(E^\times)^{p^n}=W_p.
                                                        \tag{FF.42.2}
\]
In particular no new element in a
finite extension enlarges this kernel.
The subgroup recovers its abelian
constant field by the equation
\[
                     W_p\cap(1-W_p)=F_p\setminus\{0,1\}.
                                                        \tag{FF.42.3}
\]
There are two precise value-detection
consequences.
\begin{enumerate}
\item If \(a,b\in E\setminus\{0,1\}\)
have the same pair
\((\kappa_E(a),\kappa_E(1-a))\),
then \(a=b\) or both \(a,b\in B_p\).
Thus the pair is injective outside
\(B_p\), with no collision between
such a value and a value in \(B_p\).
\item If \(a,b\in E\) satisfy
\[
                  \kappa_E(a-c)=\kappa_E(b-c)
                  \quad\text{for every }c\in\mathbb Q
                  \text{ with }a,b\ne c,
                                                        \tag{FF.42.4}
\]
then \(a=b\) or both \(a,b\in F_p\).
Conversely any two values in \(F_p\)
satisfy (FF.42.4). Thus all rational
target translates detect every value
outside \(F_p\) exactly.
\end{enumerate}
\end{claim}
\begin{proof}
Let \(a\in E^\times\) have zero
rational Kummer class. Choose a finite
Galois number field \(D/\mathbb Q\)
containing \(a\) and finitely many
generators for \(E/B_p\). We may
replace \(E\) temporarily by
\(DB_p\), since restriction preserves
the zero class. Its local completion
\(M=D_v\) is a finite Galois extension
of \(\mathbb Q_p\). With the fixed
embedding into the local algebraic
closure, put
\[
 \begin{gathered}
 I=G_{K_p},\qquad J=I\cap G_M,\qquad
 H=I''\cap J=G_{DB_p},\\
 \Gamma=J/H,\qquad C=(I'\cap J)/H.
 \end{gathered}
                                                        \tag{FF.42.5}
\]
The Galois choice of \(M\) makes
\(J\) and \(H\) normal in \(I\),
so every conjugation used below acts
on this same model. The group \(C\)
is abelian, the quotient \(\Gamma/C\)
is abelian, and the cyclotomic
character \(\chi\) is trivial on
\(C\) and has infinite image on
\(\Gamma/C\). The latter follows
because \(J\) is open in \(I\).

The class of \(a\) on \(J\)
inflates from \(H^1(\Gamma,\mathbb Q_p(1))\),
since its restriction to \(H\) is
zero. The abelian cocycle identity
used in Claim~\ref{clm:FFcyclotomiccohomology}
gives
\[
                 H^1(\Gamma/C,\mathbb Q_p(1))=0.
\]
Inflation--restriction therefore makes
the map
\[
 H^1(\Gamma,\mathbb Q_p(1))
       \longrightarrow H^1(C,\mathbb Q_p(1))
                                                        \tag{FF.42.6}
\]
injective. Every \(g\in I'\)
centralizes \(C\): the commutator
of \(g\) with an element of
\(I'\cap J\) lies both in
\(I''\) and in \(J\), hence
in \(H\). Also \(\chi(g)=1\).
Thus \(g\) acts trivially on the
right side of (FF.42.6), and hence
on its left side. Functoriality of
Kummer classes gives
\[
                         \kappa_J(g(a)/a)=0
                         \qquad(g\in I').
                                                        \tag{FF.42.7}
\]

All elements \(g(a)/a\) lie in
the same finite field \(M\).
Restriction
\(H^1(G_M,\mathbb Q_p(1))
\longrightarrow H^1(J,\mathbb Q_p(1))\)
is injective: \(J=I_M\) and
\(\mathbb Q_p(1)^J=0\), so
inflation--restriction has zero
kernel. The rational Kummer kernel
of \(M^\times\) is precisely its
finite group of roots of unity.
Indeed valuation detects the
uniformizer component, and on a
sufficiently small principal-unit
group the \(p\)-adic logarithm
identifies the remaining rationalized
Kummer completion with the additive
local field. It follows from
(FF.42.7) that
\[
                            g(a)/a\in\mu(M)
                            \qquad(g\in I').
                                                        \tag{FF.42.8}
\]
Choose a single integer \(m\)
annihilating the finite group
\(\mu(M)\). Then \(a^m\) is
fixed by every \(g\in I'\).
Since \(a\) is globally algebraic,
\(a^m\in\overline{\mathbb Q}^{I'}=F_p\).
Moreover \(\mu_m\subset F_p\),
so \(F_p(a)/F_p\) is cyclic.
The field \(B_p\) is the maximal
abelian extension of \(F_p\), with
groups \(G_{F_p}=I'\) and
\(G_{B_p}=I''\). Consequently
\(a\in B_p\), and hence
\(a\in W_p\).

Conversely suppose \(a^m\in F_p^\times\).
For every \(n\), a root \(a_n\)
with \(a_n^{p^n}=a\) satisfies
\(a_n^{mp^n}=a^m\). Because
\(F_p\) contains all roots of
unity, the extension
\(F_p(a_n)/F_p\) is cyclic and
therefore lies in \(B_p\).
Choosing these roots recursively
gives compatible divisions of \(a\)
in \(B_p\), so its integral and
rational Kummer classes vanish
over every \(E\).
Finally \(G_E\) acts trivially on
\(\mathbb Z_p(1)\), making the
integral Kummer class a homomorphism
to a torsion-free module. Its
rationalization vanishes exactly
when it vanishes integrally. This
is equivalent to divisibility by
every \(p^n\); compatible choices
exist since all \(p\)-power roots
of unity are already in \(E\).
This proves (FF.42.2).

For (FF.42.3), the inclusion of
\(F_p\setminus\{0,1\}\) is immediate.
Conversely let \(x\in W_p\) and
\(1-x\in W_p\). For every
\(g\in I'\), there are roots
of unity \(\zeta,\eta\in F_p\)
such that
\[
                    g(x)=\zeta x,\qquad
                    g(1-x)=\eta(1-x).
\]
If \(g(x)\ne x\), their sum gives
\((\zeta-\eta)x=1-\eta\).
The case \(\zeta=\eta\) would
force both to equal one, contrary
to the choice of \(g\). Otherwise
the equality puts \(x\) in \(F_p\),
also contradicting \(g(x)\ne x\).
Thus every \(g\in I'\) fixes
\(x\), and \(x\in F_p\).

For the first value-detection assertion,
equality of the two classes puts
\[
                 \alpha=a/b\in W_p,\qquad
                 \beta=(1-a)/(1-b)\in W_p.
\]
If \(\alpha=\beta\), summing the
two defining equalities gives
\(\alpha=1\), so \(a=b\).
Otherwise
\(b=(1-\beta)/(\alpha-\beta)\)
and \(a=\alpha b\) both belong
to \(B_p\).

For the second assertion assume
\(a\ne b\). Choose distinct
rational \(c_0,c_1\) avoiding
\(a,b\). The two ratios
\((a-c_i)/(b-c_i)\in W_p\subset B_p\)
are distinct; if equal, subtraction
would make both one and give
\(a=b\). Solving the two linear
equations therefore puts
\(a,b\in B_p\).
Choose a finite Galois number field
\(D\) containing them. For every
\(g\in I'\), consider
\[
 R_g(T)=\frac{(g(a)-T)(b-T)}{(g(b)-T)(a-T)}
                                 \in D(T).
                                                        \tag{FF.42.9}
\]
For each rational \(c\) away
from its finitely many exceptional
values, (FF.42.4) puts
\(\alpha_c=(a-c)/(b-c)\) in
\(W_p\). A power of \(\alpha_c\)
is fixed by \(g\), so
\(R_g(c)=g(\alpha_c)/\alpha_c\)
is a root of unity. It belongs
to the same number field \(D\).
There are only finitely many roots
of unity in \(D\), so one value
occurs for infinitely many \(c\).
A rational function with infinitely
many occurrences of one value is
constant. Its value at infinity
is one, hence \(R_g(T)=1\).
Comparing coefficients gives
\[
              g(a)-a=g(b)-b=\delta,\qquad
              \delta(b-a)=0.
\]
Since \(a\ne b\), both are
fixed by \(g\). This holds
for every \(g\in I'\), proving
\(a,b\in F_p\). The converse
follows directly from (FF.42.2).
\end{proof}

This supplies exact value detection
from the collection of all rational
target translates, despite the
failure of the single-functional
invariant lift in
Claim~\ref{clm:FFmetabelianbarrier}.
It does not yet prove reconstruction:
one must still derive sufficiently
coherent arithmetic control of those
translated Kummer maps from an
actual pinned topos map, and handle
the values in \(F_p\) which the
displayed family does not distinguish.
Nor does it make an analytic Tate
parameter globally algebraic. The
local version of the same argument
has the local fixed field
\(\overline{\mathbb Q}_p^{I'}\)
as its constants; that field already
contains \(\mathbb Q_p\) and the
local parameter \(q\) in
Claim~\ref{clm:FFmetabelianbarrier}.
